\documentclass[10pt]{article}
\usepackage[a4paper,margin=0.5in]{geometry}
\usepackage[T1]{fontenc}
\usepackage{lmodern}
\usepackage[expansion=false]{microtype}
\usepackage{tabularx,array,booktabs}
\usepackage{enumitem}
\usepackage[table]{xcolor}
\usepackage[hidelinks]{hyperref}
\usepackage[compact]{titlesec}

\titleformat*{\section}{\normalsize\bfseries}
\titlespacing*{\section}{0pt}{5pt}{2pt}
\setlength{\parindent}{0pt}
\setlength{\parskip}{3pt}
\setlength{\tabcolsep}{4pt}
\renewcommand{\arraystretch}{1.15}
\newcolumntype{L}{>{\raggedright\arraybackslash}X}
\newcolumntype{P}[1]{>{\raggedright\arraybackslash}p{#1}}
\setlist{leftmargin=*,nosep}
\definecolor{band}{gray}{0.92}
\pagestyle{empty}

\begin{document}

{\large\bfseries Fact-checking RL project: team brief\par}
\vspace{1pt}
{\small Week of 8 Oct 2026 \quad$\cdot$\quad Project ends 11 Nov 2026 \quad$\cdot$\quad Agentic lead: Chia Kwang Yang}

\section*{Background}
A fact-checker that answers confidently when the evidence is unclear spreads false confidence. Our system takes a claim and returns three things: a verdict
(Supported, Refuted or Not Enough Info), the Wikipedia sentences it relied on, and a confidence. On 900
FEVER claims, Qwen3.5-2B is right only 47\% of the time even with the correct evidence in front of it; it calls almost everything Supported.
Fine-tuning (SFT) lifts it to 78\%, above a prompted 180B model (72\% with our retrieval); it writes confidence 1.0 every time, but its verdict probability is well calibrated (error 0.03).
Our first GRPO run (no KL penalty) gained nothing: the model grew over-certain, raising that error to 0.21.
\textbf{Decided:} we train the 2B; the 4B was only about 1 point better after SFT.

\section*{What we aim to show by 11 November}
\begin{enumerate}
  \item \textbf{RL improves the verifier without breaking calibration.} The GRPO-trained 2B verifier beats the SFT model on
        accuracy or harm-weighted error (calling a false claim true is the worst mistake) while its verdict probability stays
        calibrated (error at most 0.05), on the same 900 claims, and holds up on the shortcut tests.
  \item \textbf{A cost-aware controller.} A small learned policy that decides when to search again, verify,
        abstain or answer sits above fixed policies on the accuracy-versus-search-cost curve.
  \item \textbf{Multi-turn vs single-turn RL,} same model and reward; only the loop differs.
  \item \textbf{Transfer and demo.} Results on SciFact and AVeriTeC, and a live demo that shows the decision trace.
\end{enumerate}

\section*{Who does what}
\footnotesize
\begin{tabularx}{\linewidth}{@{}P{2.95cm}P{3.6cm}LP{3.3cm}@{}}
\toprule
\textbf{Person, role} & \textbf{This week (8 Oct)} & \textbf{Weeks 2--4 (15 Oct -- 1 Nov)} & \textbf{Done when} \\
\midrule
\rowcolor{band}\multicolumn{4}{@{}l}{\textit{RL group: first GRPO run done (no gain, calibration lost). Confidence is now the verdict probability, not the written number.}} \\
\textbf{Chao Chien Hung} \newline RL-1 verifier &
  Rerun GRPO with the KL penalty on ($\beta\approx0.04$), on the 926 claims whose 8 samples score differently; stop at 200 steps if calibration slips (GPU) &
  Wk 2: final verifier, 3 seeds; fallback: SFT that writes its verdict probability (GPU). Wk 3: before/after tables. Wk 4: FEVER+SciFact vs FEVER-only run, scored on SciFact dev (GPU) &
  GRPO beats SFT on accuracy or harm, calibration error $\le$0.05; SFT vs GRPO table \\
\textbf{Michel Yildirim} \newline RL-2 reward &
  Temperature-scaling baseline on the verdict probability (the bar RL must beat); calibration metrics in AG-3's eval &
  Wk 2: 4 ablations on the fixed setup, overnight (GPU); hacking checks (confidence tracks correctness per claim, not per label). Wk 3: reliability diagrams. Wk 4: buffer &
  Reliability diagrams before/after; ablation table; hacking checks pass \\
\textbf{Yao Tianhao} \newline RL-3 controller &
  Controller env on real FEVER; abstain confidence = verifier's NEI probability &
  Wk 2: controller training (CPU). Wk 3: cost sweep; fixed-policy baselines. Wk 4: buffer &
  Accuracy-vs-cost curve; controller vs fixed policies \\
\rowcolor{band}\multicolumn{4}{@{}l}{\textit{Agentic group: the loop, the tools, the data and the demo}} \\
\textbf{Yu Jing} \newline AG-1 retrieval &
  Build the \texttt{search(query, k)} tool &
  Wk 2: sentence recall; Exa search for AVeriTeC (date cutoff, fact-check sites excluded, cached), 50-claim cost pilot. Wk 3: \texttt{follow()} hyperlink tool on FEVER; re-query &
  Retrieval recall table across datasets \\
\textbf{Chia Kwang Yang} \newline AG-2 agent loop &
  Loop skeleton over all 900 claims; one multi-turn GRPO test (GPU); keep the GPU calendar &
  Wk 2: prompted agentic baseline (GPU); loop streams each step as an event (FastAPI). Wk 3--4: multi-turn GRPO runs 1 and 2 (GPU); GRPO verifier into the demo &
  Multi-turn vs single-turn GRPO on the 900 claims; live event stream \\
\textbf{Bryan Wu Jiahe} \newline AG-3 eval \& demo &
  Eval script and frozen test splits &
  Wk 2: shortcut test set; loaders; demo Browse mode on saved traces. Wk 3: failure audit in Browse; Live mode. Wk 4: transfer runs; Compare mode &
  Every reported number comes from the eval script; demo with fallback \\
\bottomrule
\end{tabularx}
\normalsize

\section*{How we work together}
\begin{itemize}
  \item \textbf{Freeze four interfaces this week} so nobody waits on anyone: \texttt{search()} (AG-1),
        \texttt{verify()} (RL-1 trains it, AG-2 sets the signature), \texttt{reward()} (RL-2), and the
        trace JSONL, one row per claim (AG-2 writes it; RL-3 trains on it; AG-3 scores it).
  \item \textbf{No number goes on a slide unless it came from AG-3's eval script.}
  \item \textbf{One GPU} (RTX 5070 Ti, 16GB). Week 1: RL-1 by day; one AG-2 test. Week 2: RL-1 final verifier by day; AG-2 baseline and RL-2
        ablations share nights. Week 3: AG-2 first, leftover ablations when free. Week 4: AG-2 by day, RL-1 SciFact run overnight.
        Week 5: spare. Book with KY.
  \item \textbf{Reviews:} 15 Oct (does fixed GRPO beat SFT and keep calibration?), 29 Oct (pick the lead claim),
        5--11 Nov (report, slides, demo dry run, final review).
  \item \textbf{Demo:} one frontend fed by the loop's trace events. \emph{Live}: any claim, steps appear as they run (Exa for open claims).
        \emph{Compare}: SFT vs GRPO on the same claim. \emph{Browse}: saved traces; doubles as the failure audit and the fallback.
  \item \textbf{Tools by dataset:} FEVER uses Wikipedia search plus hyperlink follow; SciFact searches its abstracts; AVeriTeC and the demo use Exa.
  \item \textbf{Open:} AVeriTeC's fourth label (Conflicting/Cherrypicking) mapped to NEI or dropped, AG-3 to decide. Exa key: KY.
        Out of scope: dense retriever, escalate action, LIAR-PLUS.
\end{itemize}

\end{document}
